表\ref{tab:obsindex}に各クラスの指標値（中央値）を示す。高被引用論文ではクラスG（成長論文）が最も多く、次にクラスNm、クラスOと続き、クラスFPとクラスSBは同程度に出現している。なお、クラスFP（Flash in the pan）の年齢が高いのは、このクラスに分類されるには、早期の高い被引用だけでなくその後の十分な低い被引用期間も要することによる。
表\ref{tab:mcpobsindex}に最多被引用論文の指標値を示す。最多被引用論文の各指標値の範囲は広く、必ずしも特定のクラスから出現するわけではないことを示唆している。なお、うち2本はクラスOに分類されているが、被引用曲線を確認したところ明確な多峰性が見られ、リバイバル論文であることが明らかであった。
